% !TEX root = main.tex
\chapter{Εισαγωγή}
\label{chap:introduction}

\section{Το Πρόβλημα Δειγματοληψίας σε Υψηλές Διαστάσεις}
\label{sec:intro-sampling}

Ένα κεντρικό ζήτημα στην υπολογιστική στατιστική, τη μηχανική μάθηση και τη
στατιστική φυσική είναι η παραγωγή δειγμάτων από μια κατανομή-στόχο πιθανότητας
$\pi$ στο $\mathbb{R}^d$ που είναι γνωστή μόνο έως μια σταθερά κανονικοποίησης.
Η συνήθης διατύπωση είναι να εκφραστεί η κατανομή-στόχος ως \emph{μέτρο Gibbs}
\begin{equation}
  \pi(x) \;\propto\; \exp(-U(x)), \qquad x \in \mathbb{R}^d,
  \label{eq:gibbs-intro}
\end{equation}
όπου $U : \mathbb{R}^d \to \mathbb{R}$ είναι μια συνάρτηση δυναμικού. 
\newline
Η άγνωστη σταθερά κανονικοποίησης $\int_{\mathbb{R}^d} \exp(-U(y))\, dy$ μπορεί να είναι
αδύνατο να υπολογιστεί (ιδιαίτερα σε υψηλές διαστάσεις), ωστόσο ένα ευρύ φάσμα ερωτημάτων στατιστικής
συμπερασματολογίας μπορεί να απαντηθεί αποκλειστικά από δείγματα: οι εκ των
υστέρων κατανομές στην Μπεϋζιανή στατιστική, οι ελεύθερες ενέργειες στη
στατιστική μηχανική, οι περιθώριες πιθανοφάνειες στην επιλογή μοντέλου ανάγονται όλα στον υπολογισμό του $\mathbb{E}_\pi[\varphi(X)]$ για κάποια συνάρτηση ελέγχου $\varphi$.

Σε χαμηλή διάσταση, γενικές μέθοδοι όπως το rejection ή το importance sampling λειτουργούν καλά. Καθώς η διάσταση αυξάνεται, ωστόσο,
αυτές οι μέθοδοι αποτυγχάνουν: ο όγκος του $\mathbb{R}^d$ εκρήγνυται ταχύτερα
από ό,τι μπορεί να καλύψει οποιοδήποτε σταθερό περίβλημα ή κατανομή πρότασης,
οπότε τα περισσότερα προτεινόμενα δείγματα είτε απορρίπτονται (στο rejection sampling) είτε λαμβάνουν βάρη σχεδόν μηδενικά (στο importance sampling), και η διακύμανση του εκτιμητή καθίσταται καταστροφική
\citep{robert2004monte}.

Η πιο διαδεδομένη προσέγγιση είναι το \emph{Markov chain Monte Carlo} (MCMC). Αντί να
παράγουμε ανεξάρτητα δείγματα, παράγουμε μια ενιαία ακολουθία
$(X_k)_{k\geq 0}$ \emph{εξαρτημένων} δειγμάτων, η οποία κατανέμεται
ασυμπτωτικά σύμφωνα με την $\pi$. Συγκεκριμένα, σχεδιάζουμε μια
αλυσίδα Markov της οποίας η μοναδική στάσιμη κατανομή είναι η $\pi$. Η εκτέλεσή
της για αρκετά μεγάλο χρόνο παράγει δείγματα που (αν και μη ανεξάρτητα)
συμπεριφέρονται ως τυχαίες επιλογές από την $\pi$ για τον σκοπό του υπολογισμού
προσδοκιών. Η μαθηματική εγγύηση που δικαιολογεί αυτό είναι το \emph{εργοδικό
θεώρημα}: υπό ήπιες συνθήκες, ο χρονικός μέσος $n^{-1}\sum_{k=1}^n \varphi(X_k)$
συγκλίνει στην προσδοκία $\mathbb{E}_\pi[\varphi(X)]$ καθώς $n \to \infty$
\citep{meyn2009markov, tierney1994markov}. Ο ρυθμός σύγκλισης εξαρτάται από τον
βαθμό συσχέτισης διαδοχικών δειγμάτων. Το καθιερωμένο μέτρο αυτού είναι το
\emph{αποτελεσματικό μέγεθος δείγματος} (ESS), που μετρά τον αριθμό των
πραγματικά ανεξάρτητων δειγμάτων σε μια συσχετισμένη αλυσίδα.

Η οικογένεια αλγορίθμων Langevin MCMC είναι από τις πιο εκτενώς μελετημένες και
ευρέως χρησιμοποιούμενες τεχνικές δειγματοληψίας \citep{roberts1996exponential, dalalyan2017theoretical, durmus2017nonasymptotic}. Κάθε αλγόριθμος αυτής της οικογένειας προκύπτει από τη
διακριτοποίηση μιας συνεχούς-χρόνου στοχαστικής διαφορικής εξίσωσης (SDE) της
οποίας το αναλλοίωτο μέτρο είναι ακριβώς το $\pi$, έτσι ώστε η δειγματοληψία από
το $\pi$ να ανάγεται σε πρόβλημα αριθμητικής ολοκλήρωσης μιας διάχυσης. Η διάχυση
Langevin έχει το ελκυστικό χαρακτηριστικό ότι χρησιμοποιεί μόνο την κλίση
$\nabla U$ της λογαριθμικής πυκνότητας της κατανομής-στόχου --- αποφεύγοντας τη
σταθερά κανονικοποίησης --- και κληρονομεί εργοδικότητα από την υποκείμενη στοχαστική διαφορική υπό ήπιες συνθήκες κανονικότητας του $U$ \citep{roberts1996exponential,
pavliotis2014stochastic}.

\section{Η Δυναμική Langevin ως Εργαλείο Δειγματοληψίας}
\label{sec:intro-langevin}

Η overdamped SDE Langevin στο $\mathbb{R}^d$,
\begin{equation}
  dX_t = -\nabla U(X_t)\, dt + \sqrt{2}\, dW_t,
  \label{eq:overdamped-intro}
\end{equation}
έχει ως μοναδική αναλλοίωτη κατανομή το μέτρο Gibbs $\pi \propto e^{-U}$ και
συγκλίνει στην $\pi$ εκθετικά γρήγορα όταν η $U$ είναι ισχυρά κυρτή (strongly convex) \citep{roberts1996exponential, bakry2014analysis}. Η διακριτοποίηση
του~\eqref{eq:overdamped-intro} με βήμα Euler--Maruyama μεγέθους $\eta > 0$ οδηγεί
στον \emph{Unadjusted Langevin Algorithm} (ULA), τον απλούστερο αλγόριθμο αυτής
της οικογένειας.  Ο ULA όμως είναι ασυμπτωτικά μεροληπτικός: λόγω του
  σφάλματος διακριτοποίησης, συγκλίνει σε λανθασμένη στάσιμη κατανομή
  και όχι στην $\pi$, ανεξαρτήτως πόσα δείγματα συλλέγονται
  \citep{roberts2001ula}. Η διόρθωση του σφάλματος διακριτοποίησης με βήμα
Metropolis--Hastings οδηγεί στον \emph{Metropolis-Adjusted Langevin Algorithm}
(MALA), ο οποίος δειγματολαμβάνει ακριβώς από την $\pi$ με κόστος ένα επιπλέον
έλεγχο αποδοχής/απόρριψης \citep{roberts1998optimal, roberts2001ula}.

Μια εναλλακτική δυναμική διατύπωση επαυξάνει τη θέση $X$ με μια μεταβλητή
ορμής $V$ που υπακούει σε μια \emph{underdamped} SDE Langevin,
\begin{equation}
  dX_t = V_t\, dt, \qquad
  dV_t = -\nabla U(X_t)\, dt - \gamma V_t\, dt + \sqrt{2\gamma}\, dW_t,
  \label{eq:underdamped-intro}
\end{equation}
η οποία αντιστοιχεί στο kinetic καθεστώς της στατιστικής μηχανικής
\citep{leimkuhler2015molecular}. Η ορμή επιτρέπει στην αλυσίδα να διατηρεί
πληροφορία κατεύθυνσης κατά τις επαναλήψεις, δίνοντας βελτιωμένη ανάμειξη σε
πολλά καθεστώτα. Αριθμητικά σχήματα τύπου splitting  όπως το σχήμα BAOAB
\citep{leimkuhler2013bao} παρέχουν ακρίβεια υψηλής τάξης σε τετραγωνικά δυναμικά
μαζί με ελεγχόμενη μεροληψία διακριτοποίησης για μη τετραγωνικό $U$.

Αυτοί οι τρεις αλγόριθμοι --- ULA, MALA, BAOAB --- αντιπροσωπεύουν τρία σημεία
σε ένα φάσμα συμβιβασμών ακρίβειας--κόστους. Ο ULA είναι ο φθηνότερος ανά
επανάληψη αλλά έχει μόνιμη μεροληψία $\mathcal{O}(\eta)$ \citep{durmus2017nonasymptotic}.
Ο MALA είναι ακριβής για οποιοδήποτε $\eta$ αλλά πληρώνει το κόστος μιας
διόρθωσης Metropolis και εκθέτει μια ποινή αποδοχής ανάλογη της διάστασης
\citep{roberts1998optimal}. Ο BAOAB είναι μεροληπτικός όπως ο ULA αλλά ωφελείται
από τα δομικά πλεονεκτήματα της kinetic διατύπωσης.

\section{Ο Ρόλος των Ουρών στην δειγματοληψία}
\label{sec:intro-tails}

Είναι γνωστό από τη βιβλιογραφία ότι η απόδοση των αλγορίθμων MCMC εξαρτάται
σε μεγάλο βαθμό από τη συμπεριφορά ουράς της κατανομής-στόχου \citep{jarner2007convergence}. Για κατανομές-στόχους με ελαφριές ουρές --- εκείνες για τις
οποίες η $\pi$ φθίνει τουλάχιστον τόσο γρήγορα όσο μια Γκαουσιανή μακριά από
την κορυφή του --- η γεννήτρια συνάρτηση ροπών υπάρχει και εφαρμόζεται η τυπική
θεωρία σύγκλισης. Για κατανομές-στόχους με \emph{βαριές ουρές}, όπου η γεννήτρια
συνάρτηση ροπών δεν υπάρχει σε κανένα θετικό διάστημα και το δυναμικό $U$
αυξάνεται υποτετραγωνικά (ή καθόλου) μακριά από την κορυφή, η κατάσταση είναι
εμφανώς διαφορετική.

Ο μηχανισμός είναι γεωμετρικός. Η ντετερμινιστική ολίσθηση της overdamped SDE
Langevin τραβά την αλυσίδα πίσω προς την κορυφή με ταχύτητα ανάλογη του
$\nabla U$. Για ένα Γκαουσιανό δυναμικό $U(x) = \tfrac12 \|x\|^2$, αυτή η
επαναφορική δύναμη αυξάνεται γραμμικά με το $\|x\|$, οπότε η αλυσίδα επιστρέφει
γρήγορα από οποιαδήποτε εκδρομή στην ουρά. Για ένα δυναμικό Cauchy
$U(x) = \tfrac{d+1}{2}\log(1 + \|x\|^2)$, η επαναφορική δύναμη έχει αντίθετη
ασυμπτωτική συμπεριφορά: το $\nabla U(x) = (d+1) x / (1 + \|x\|^2)$ τείνει στο
\emph{μηδέν} καθώς $\|x\|\to\infty$. Μόλις η αλυσίδα φτάσει σε ένα αρκετά
απομακρυσμένο σημείο, δεν υφίσταται ουσιαστικά καμία ολίσθηση, και ο μόνος
μηχανισμός που την εμποδίζει να απομακρυνθεί περαιτέρω είναι ο όρος θορύβου
\citep{jarner2007convergence}.

Αυτή η έλλειψη επαναφορικής δύναμης δημιουργεί αρκετά πρακτικά προβλήματα.
Για αρχή, ο αλγόριθμος δυσκολεύεται να φτάσει τις βαθιές ουρές: το αναμενόμενο
μέγεθος κάθε βήματος δεν ξεπερνά το $\sqrt{2\eta}$, οπότε η αλυσίδα χρειάζεται
πολλά βήματα για να φτάσει μακριά, ενώ χωρίς ολίσθηση τέτοιες εκδρομές γίνονται
σπάνιες. Επίσης, όταν συμβεί μια εκδρομή, η αλυσίδα τείνει να κολλάει στην ουρά
για πολλές επαναλήψεις, αυξάνοντας τον χρόνο αυτοσυσχέτισης. Για τον MALA
συγκεκριμένα, τα μεγάλα άλματα \emph{προς} τη βαριά ουρά γίνονται εύκολα
αποδεκτά επειδή η λογαριθμική πυκνότητα είναι σχεδόν επίπεδη εκεί, αλλά τα
άλματα \emph{από} την ουρά πίσω στο κεντρικό σώμα μπορεί να απορριφθούν λόγω
ασυμμετρίας της κατανομής πρότασης --- και αυτό επιδεινώνεται σε υψηλότερη
διάσταση. Το αποτέλεσμα είναι ότι κανένας αλγόριθμος της οικογένειας Langevin
δεν δειγματοληπτεί αξιόπιστα κατανομές βαριών ουρών χωρίς προσεκτική ρύθμιση.

\section{Κεντρική Υπόθεση}
\label{sec:intro-hypothesis}

Με βάση τα παραπάνω και την υπάρχουσα βιβλιογραφία για τους δειγματολήπτες
Langevin, η μελέτη αυτή οργανώνεται γύρω από την ακόλουθη υπόθεση.

\begin{quote}
\textbf{Υπόθεση.}
Ο ULA έχει μόνιμη μεροληψία λόγω σφάλματος διακριτοποίησης
\citep{durmus2017nonasymptotic}, ενώ η διόρθωση Metropolis--Hastings στον MALA
την εξαλείφει \citep{roberts1998optimal}. Οι κατανομές με βαριές ουρές
δυσκολεύουν τις overdamped παραλλαγές γιατί η επαναφορική δύναμη της SDE
Langevin εξασθενεί μακριά από την κορυφή. Ο ULA επηρεάζεται περισσότερο, ενώ
ο MALA αντισταθμίζει εν μέρει επιλέγοντας μεγαλύτερο βήμα. Το kinetic σχήμα
(BAOAB) δεν χρειάζεται βήμα Metropolis για να ελέγξει την μεροληψία
\citep{leimkuhler2013bao} και αναμένεται να αποδίδει καλύτερα στα ακραία
ποσοστημόρια κατανομών με βαριές ουρές, εκεί όπου ο MALA δυσκολεύεται με την
ασύμμετρη γεωμετρία πρότασης-αποδοχής. Τέλος, η επιλογή της αριθμητικής
διακριτοποίησης (Euler--Maruyama ή BAOAB) αναμένεται να έχει σημαντική επίδραση
στην ποιότητα δειγματοληψίας, ανεξάρτητα από το αν υπάρχει διόρθωση Metropolis.
\end{quote}

Η εργασία ελέγχει εμπειρικά αυτή την υπόθεση σε τρεις
κατανομές-στόχους που καλύπτουν ένα ελεγχόμενο εύρος συμπεριφοράς ουράς --- μια
Γκαουσιανή (ελαφριές ουρές), μια Student-$t$ με $\nu=5$ βαθμούς ελευθερίας
(πολυωνυμικές ουρές, πεπερασμένη διακύμανση) και μια τυπική κατανομή Cauchy
(πολυωνυμικές ουρές, χωρίς πεπερασμένες ροπές) --- σε διαστάσεις από $d=1$ έως
$d=100$, με ένα ποσοτικό σύνολο διαγνωστικών που περιλαμβάνει αποτελεσματικό
μέγεθος δείγματος, αυτοσυσχέτιση, σφάλματα ποσοστημορίων, κάλυψη ουράς και
μέτρα καλής προσαρμογής.

\section{Δομή της Εργασίας}
\label{sec:intro-outline}

Το υπόλοιπο της εργασίας οργανώνεται ως εξής.
Το Κεφάλαιο~\ref{chap:background} παρουσιάζει το θεωρητικό και μαθηματικό
υπόβαθρο: στοχαστικές διαδικασίες, θεωρία αλυσίδων Markov και το εργοδικό
θεώρημα, η SDE Langevin και το αναλλοίωτο μέτρο της, και ο χαρακτηρισμός
κατανομών με βαριές ουρές μέσω της αποτυχίας των εκθετικών ροπών.
Το Κεφάλαιο~\ref{chap:algorithms} περιγράφει τους τρεις αλγορίθμους βασισμένους
στο Langevin που μελετώνται στην εργασία: τους overdamped ULA και MALA, και τον
underdamped BAOAB splitting ολοκληρωτή, μαζί με συζήτηση του ρόλου της ορμής και
της επιλογής σχήματος διακριτοποίησης.
Το Κεφάλαιο~\ref{chap:experiments} περιγράφει την πειραματική ρύθμιση και
παρουσιάζει τα κύρια αποτελέσματα από τη μελέτη πλέγματος, τα διαγνωστικά
μελετών περίπτωσης, την ανάλυση συσχέτισης βαρύτητας ουράς και τη μελέτη
σύγκλισης.
Το Κεφάλαιο~\ref{chap:conclusions} συνθέτει τα ευρήματα σε συστάσεις ανάπτυξης,
συζητά περιορισμούς και επιφυλάξεις, και σκιαγραφεί κατευθύνσεις για μελλοντική
έρευνα. Το Παράρτημα περιέχει ψευδοκώδικα αλγορίθμων, υποστηρικτικούς πίνακες
αποτελεσμάτων και επιπλέον διαγνωστικά διαγράμματα.
